\documentclass[11pt,a4paper]{article}

% ================ Encodage et langue ================
\usepackage[utf8]{inputenc}
\usepackage[T1]{fontenc}
\usepackage[french]{babel}
\usepackage{lmodern}
\usepackage{textcomp}

% ================ Mise en page ================
\usepackage[margin=2.2cm]{geometry}
\usepackage{parskip}
\setlength{\parindent}{0pt}
\setlength{\parskip}{0.6em}

% ================ Couleurs et liens ================
\usepackage[table]{xcolor}
\definecolor{primary}{HTML}{0B5394}
\definecolor{accent}{HTML}{1F77B4}
\definecolor{light}{HTML}{F2F2F2}
\definecolor{codebg}{HTML}{F6F8FA}
\definecolor{rulec}{HTML}{D0D7DE}

\usepackage[colorlinks=true, linkcolor=primary, urlcolor=primary, citecolor=primary]{hyperref}

% ================ Tableaux ================
\usepackage{longtable, booktabs, array, tabularx}
\renewcommand{\arraystretch}{1.25}
\arrayrulecolor{rulec}

% ================ Listes ================
\usepackage{enumitem}
\setlist{itemsep=2pt, topsep=4pt}

% ================ Code et citations ================
\usepackage{fancyvrb}
\usepackage{framed}
\definecolor{shadecolor}{named}{light}

% Encadrés pour les blocs > (quotes)
\usepackage{mdframed}
\newmdenv[
  linecolor=primary,
  linewidth=2pt,
  topline=false, bottomline=false, rightline=false,
  innerleftmargin=10pt, innerrightmargin=8pt,
  innertopmargin=6pt, innerbottommargin=6pt,
  backgroundcolor=codebg,
  skipabove=8pt, skipbelow=8pt
]{quotebox}
\renewenvironment{quote}{\begin{quotebox}}{\end{quotebox}}

% ================ Titres ================
\usepackage{titlesec}
\setcounter{secnumdepth}{0}
\titleformat{\section}{\Large\bfseries\color{primary}}{\thesection.}{0.6em}{}
\titleformat{\subsection}{\large\bfseries\color{accent}}{\thesubsection}{0.6em}{}
\titleformat{\subsubsection}{\normalsize\bfseries}{\thesubsubsection}{0.6em}{}

% ================ En-têtes et pieds de page ================
\usepackage{fancyhdr}
\pagestyle{fancy}
\fancyhf{}
\fancyhead[L]{\small\color{gray}ESPRIT School of Business --- L3 Business Computing BI --- Projet Intégré}
\fancyfoot[C]{\small\thepage}
\renewcommand{\headrulewidth}{0.3pt}
\renewcommand{\footrulewidth}{0pt}

% ================ Pandoc helpers ================
\providecommand{\tightlist}{\setlength{\itemsep}{0pt}\setlength{\parskip}{0pt}}
\providecommand{\passthrough}[1]{#1}

\usepackage{calc}
\providecommand{\real}[1]{#1}


% ================ Métadonnées ================
\title{\color{primary}\textbf{Projet Intégré --- Analyse et Prédiction
du Churn Client}}
\author{Aymen Ben Brik et Hela Hammami}
\date{2026}

\begin{document}

\begin{titlepage}
\centering
\vspace*{1.5cm}

{\Large\bfseries ESPRIT School of Business}\\[0.3cm]
{\large Licence 3 --- Business Computing BI}\\[2cm]

{\color{gray}\large Projet Intégré}\\[0.3cm]
{\color{gray}\large Analyse et Prédiction du Churn Client}\\[2.2cm]

\rule{0.85\linewidth}{1.2pt}\\[0.6cm]
{\Huge\bfseries\color{primary} Projet Intégré --- Analyse et Prédiction
du Churn Client\par}
\rule{0.85\linewidth}{1.2pt}\\[2cm]

{\large Tuteurs du projet}\\[0.2cm]
{\Large\bfseries Aymen Ben Brik et Hela Hammami}\\[1cm]
{\large 2026}

\vfill
{\color{gray}\small Document à destination des étudiants}
\end{titlepage}

\tableofcontents
\newpage

\section{1. Contexte}\label{contexte}

Un établissement financier met à votre disposition un jeu de données
anonymisé et généralisé décrivant des clients, leurs comptes et leurs
produits. La fidélisation client (lutte contre l'attrition, ou
\emph{churn}) est un enjeu majeur du secteur : acquérir un nouveau
client coûte plusieurs fois plus cher que d'en conserver un existant.

Votre mission est de concevoir, de bout en bout, une \textbf{solution
décisionnelle complète} permettant de comprendre, mesurer et anticiper
le churn, puis de mettre ces analyses à disposition des décideurs métier
via une interface web.

\section{2. Objectifs pédagogiques}\label{objectifs-puxe9dagogiques}

Ce projet vous permet de mettre en pratique l'ensemble de la chaîne
décisionnelle :

\begin{itemize}
\tightlist
\item
  Comprendre et qualifier des données (analyse univariée, valeurs
  manquantes, outliers).
\item
  Concevoir un entrepôt de données et ses flux d'alimentation
  (modélisation dimensionnelle, ETL).
\item
  Produire des indicateurs métier et des tableaux de bord avec Power BI
  et DAX.
\item
  Modéliser l'évolution d'un KPI dans le temps (série temporelle, Deep
  Learning) et prédire le churn (Machine Learning).
\item
  Concevoir et livrer une interface web documentée par des diagrammes
  UML.
\item
  Travailler en équipe avec une démarche projet structurée (versioning,
  livrables, présentation).
\end{itemize}

\section{3. Périmètre fonctionnel}\label{puxe9rimuxe8tre-fonctionnel}

Le projet est structuré en \textbf{cinq phases successives}. Chaque
phase comporte un socle \textbf{obligatoire} et des \textbf{bonus}
optionnels. Traitez toujours le socle obligatoire avant d'aborder un
bonus.

{\def\LTcaptype{none} % do not increment counter
\begin{longtable}[]{@{}
  >{\raggedright\arraybackslash}p{(\linewidth - 8\tabcolsep) * \real{0.1176}}
  >{\raggedright\arraybackslash}p{(\linewidth - 8\tabcolsep) * \real{0.1765}}
  >{\raggedright\arraybackslash}p{(\linewidth - 8\tabcolsep) * \real{0.1176}}
  >{\raggedright\arraybackslash}p{(\linewidth - 8\tabcolsep) * \real{0.2941}}
  >{\raggedright\arraybackslash}p{(\linewidth - 8\tabcolsep) * \real{0.2941}}@{}}
\toprule\noalign{}
\begin{minipage}[b]{\linewidth}\raggedright
Phase
\end{minipage} & \begin{minipage}[b]{\linewidth}\raggedright
Intitulé
\end{minipage} & \begin{minipage}[b]{\linewidth}\raggedright
Échéance
\end{minipage} & \begin{minipage}[b]{\linewidth}\raggedright
Obligatoire
\end{minipage} & \begin{minipage}[b]{\linewidth}\raggedright
Bonus
\end{minipage} \\
\midrule\noalign{}
\endhead
\bottomrule\noalign{}
\endlastfoot
1 & Collecte et compréhension des données & Semaine du 9/11 & Analyse
univariée ; décompte des outliers et des valeurs manquantes & Web
scraping (Beautiful Soup, Playwright, Scrapy) ; analyse bivariée \\
2 & Conception du DWH et ETL & Semaine du 16/11 & Talend + PostgreSQL
administré avec pgAdmin & SSIS ; Databricks (bronze / silver / gold) ;
SCD de type 2 ; Python \\
3 & Reporting et dashboard & 23/11 & Power BI ; DAX pour les KPI &
Visualisation avec Databricks ; interface JavaScript \\
4 & Modèles prédictifs & 1/12 & Fixer un KPI ; analyse de sa série
temporelle ; modèle de série temporelle en Deep Learning (LSTM / RNN) ;
Machine Learning pour le churn & Traiter plusieurs KPI \\
5 & Déploiement & 7/12 & Interface web ; diagrammes UML : cas
d'utilisation global, une séquence, classes & DevOps : tests ;
application mobile \\
\end{longtable}
}

\subsection{3.1 Phase 1 --- Collecte et compréhension des
données}\label{phase-1-collecte-et-compruxe9hension-des-donnuxe9es}

\begin{itemize}
\tightlist
\item
  Charger les fichiers fournis et vérifier leur intégrité.
\item
  Réaliser l'\textbf{analyse univariée} de chaque variable utile
  (distribution, tendance centrale, dispersion, modalités).
\item
  Produire le \textbf{décompte des valeurs manquantes} et le
  \textbf{décompte des outliers} par variable, en précisant la méthode
  retenue (IQR, z-score, etc.).
\item
  \emph{Bonus} : enrichir le contexte par \textbf{web scraping} ;
  réaliser une \textbf{analyse bivariée}.
\end{itemize}

\subsection{3.2 Phase 2 --- Conception du DWH et
ETL}\label{phase-2-conception-du-dwh-et-etl}

\begin{itemize}
\tightlist
\item
  Concevoir un modèle dimensionnel en étoile (faits + dimensions) et
  justifier le grain.
\item
  Implémenter les flux d'alimentation avec \textbf{Talend}.
\item
  Créer et administrer l'entrepôt sous \textbf{PostgreSQL} avec
  \textbf{pgAdmin}.
\item
  \emph{Bonus} : version \textbf{SSIS} ; architecture
  \textbf{Databricks} en médaillon (bronze / silver / gold) ;
  historisation \textbf{SCD de type 2} ; traitements en \textbf{Python}.
\end{itemize}

\subsection{3.3 Phase 3 --- Reporting et
dashboard}\label{phase-3-reporting-et-dashboard}

\begin{itemize}
\tightlist
\item
  Produire un rapport \textbf{Power BI} connecté à l'entrepôt.
\item
  Écrire les KPI en \textbf{DAX} et permettre une exploration
  interactive (filtres, drill-down).
\item
  Au minimum : vue d'ensemble, analyse du churn, segmentation client.
\item
  \emph{Bonus} : visualisation avec \textbf{Databricks} ; restitution
  dans une \textbf{interface JavaScript}.
\end{itemize}

\subsection{3.4 Phase 4 --- Modèles
prédictifs}\label{phase-4-moduxe8les-pruxe9dictifs}

\begin{itemize}
\tightlist
\item
  \textbf{Fixer un KPI} et justifier son intérêt métier (exemple :
  nombre hebdomadaire de clôtures de comptes).
\item
  Analyser sa \textbf{série temporelle} : tendance, saisonnalité,
  ruptures, pics.
\item
  Construire un modèle de série temporelle en \textbf{Deep Learning
  (LSTM ou RNN)}, comparé à une référence simple.
\item
  Construire un modèle de \textbf{Machine Learning} de prédiction du
  churn et l'évaluer avec des métriques adaptées.
\item
  \emph{Bonus} : appliquer la démarche à \textbf{plusieurs KPI}.
\end{itemize}

\subsection{3.5 Phase 5 --- Déploiement}\label{phase-5-duxe9ploiement}

\begin{itemize}
\tightlist
\item
  Développer une \textbf{interface web} exposant les prédictions et les
  indicateurs de synthèse.
\item
  Livrer trois \textbf{diagrammes UML} : cas d'utilisation
  \textbf{global}, \textbf{une seule} séquence, classes.
\item
  \emph{Bonus} : démarche \textbf{DevOps} avec tests automatisés ;
  \textbf{application mobile}.
\end{itemize}

\section{4. Livrables attendus}\label{livrables-attendus}

Trois livrables sont attendus à la fin du projet, en plus du rendu de
chaque phase :

{\def\LTcaptype{none} % do not increment counter
\begin{longtable}[]{@{}
  >{\raggedright\arraybackslash}p{(\linewidth - 4\tabcolsep) * \real{0.3333}}
  >{\raggedright\arraybackslash}p{(\linewidth - 4\tabcolsep) * \real{0.3333}}
  >{\raggedright\arraybackslash}p{(\linewidth - 4\tabcolsep) * \real{0.3333}}@{}}
\toprule\noalign{}
\begin{minipage}[b]{\linewidth}\raggedright
Livrable
\end{minipage} & \begin{minipage}[b]{\linewidth}\raggedright
Format
\end{minipage} & \begin{minipage}[b]{\linewidth}\raggedright
Contenu attendu
\end{minipage} \\
\midrule\noalign{}
\endhead
\bottomrule\noalign{}
\endlastfoot
\textbf{Code source} & Dépôt GitHub, sans aucune donnée & Notebooks,
jobs Talend exportés, scripts SQL, mesures DAX, code ML, application
web, diagrammes UML, README détaillé, instructions d'installation et
d'exécution \\
\textbf{Présentation orale} & Slides (PPT/PDF), 20 min + 10 min de
questions & Démonstration de l'application, présentation de la démarche,
principaux résultats \\
\textbf{Rapport} & PDF, 25--40 pages & Document professionnel décrivant
la démarche, l'architecture, les choix techniques, les résultats, les
limites et les perspectives \\
\end{longtable}
}

\section{5. Modalités}\label{modalituxe9s}

\begin{itemize}
\tightlist
\item
  \textbf{Composition des équipes} : 3 à 4 étudiants par groupe.
\item
  \textbf{Durée} : cinq jalons, de la semaine du 9/11 au 7/12/2026, à
  temps partagé avec les autres cours.
\item
  \textbf{Encadrement} : un point d'avancement à chaque jalon avec les
  tuteurs.
\item
  \textbf{Outils} : les outils marqués \emph{obligatoire} sont imposés.
  Les autres restent libres dans le respect du guide étudiant (document
  4).
\end{itemize}

\section{6. Critères d'évaluation}\label{crituxe8res-duxe9valuation}

{\def\LTcaptype{none} % do not increment counter
\begin{longtable}[]{@{}
  >{\raggedright\arraybackslash}p{(\linewidth - 2\tabcolsep) * \real{0.5000}}
  >{\raggedright\arraybackslash}p{(\linewidth - 2\tabcolsep) * \real{0.5000}}@{}}
\toprule\noalign{}
\begin{minipage}[b]{\linewidth}\raggedright
Critère
\end{minipage} & \begin{minipage}[b]{\linewidth}\raggedright
Pondération
\end{minipage} \\
\midrule\noalign{}
\endhead
\bottomrule\noalign{}
\endlastfoot
Compréhension des données (analyse univariée, manquants, outliers) & 10
\% \\
Qualité du DWH et de l'ETL (Talend, PostgreSQL) & 20 \% \\
Pertinence du reporting (Power BI, DAX) & 15 \% \\
Rigueur des modèles prédictifs (série temporelle, churn) & 25 \% \\
Interface web et diagrammes UML & 15 \% \\
Qualité du code, du rapport et de la présentation orale & 15 \% \\
\end{longtable}
}

Les \textbf{bonus} sont évalués en plus du socle et ne compensent pas
une partie obligatoire absente.

\section{7. Règles importantes}\label{ruxe8gles-importantes}

\begin{itemize}
\tightlist
\item
  \textbf{Confidentialité} : les données sont confidentielles et remises
  hors de GitHub. Aucune donnée, même partielle ou agrégée à la ligne,
  ne doit être publiée (dépôt, application, capture d'écran, modèle
  embarquant des lignes). Seul le code peut être partagé.
\item
  \textbf{Non-réidentification} : toute tentative d'identifier
  l'établissement, une agence ou un client est \textbf{interdite}, y
  compris par recoupement ou web scraping. Si vous pensez reconnaître
  l'une de ces informations, n'en parlez pas dans vos livrables et
  prévenez immédiatement les tuteurs.
\item
  \textbf{Anonymat} : ne formulez aucune hypothèse sur l'origine des
  données, ni dans le code, ni dans le rapport, ni dans la présentation.
\item
  \textbf{Application publique} : l'interface déployée n'affiche que des
  agrégats et des profils fictifs. Toute liste de clients reste derrière
  une authentification et n'est jamais incluse dans le dépôt.
\item
  \textbf{Versioning} : le dépôt GitHub doit montrer un historique
  régulier de commits par chaque membre de l'équipe.
\item
  \textbf{Reproductibilité} : tout doit pouvoir être ré-exécuté par un
  tiers à partir du README.
\end{itemize}

\end{document}
